\exercise{14.1}{Six initial convergence tests}
\goal Decide convergence for each of the six given series, with a
reason that checks the hypotheses of the chosen test.

\step{(a) The series with term $n^4/2^n$.}
Let $a_n=n^4/2^n>0$. Since no term is zero, the successive ratio is
well defined. Canceling powers gives
\[
 \frac{a_{n+1}}{a_n}
 =\frac{(n+1)^4}{2^{n+1}}\frac{2^n}{n^4}
 =\frac12\left(1+\frac1n\right)^4
 \longrightarrow\frac12.
\]
The last limit follows from $1/n\to0$ and the finite product limit
law. Its value is strictly below $1$, so the Ratio Test proves
absolute convergence. Since the terms are positive, this proves
convergence of the requested series.

\step{(b) The series with term $2^n/n!$.}
Put $a_n=2^n/n!>0$. The identity $(n+1)!=(n+1)n!$ yields
\[
 \frac{a_{n+1}}{a_n}
 =\frac{2^{n+1}}{(n+1)!}\frac{n!}{2^n}
 =\frac{2}{n+1}\longrightarrow0.
\]
Here $0<1$, so the Ratio Test proves absolute convergence.

\step{(c) The series with term $n^2/3^n$.}
With $a_n=n^2/3^n>0$, the same cancellation gives
\[
 \frac{a_{n+1}}{a_n}
 =\frac13\left(1+\frac1n\right)^2\longrightarrow\frac13<1.
\]
The Ratio Test again gives absolute convergence. The power of $n$
in the numerator does not prevent this conclusion because the
computed ratio limit is strictly less than $1$.

\step{(d) The series with term $n!/(n^4+3)$.}
Let $a_n=n!/(n^4+3)>0$. Its ratio is
\[
 \frac{a_{n+1}}{a_n}
 =(n+1)\frac{n^4+3}{(n+1)^4+3}
 =(n+1)\frac{1+3/n^4}{(1+1/n)^4+3/n^4}.
\]
The second factor tends to $1$, since its numerator and denominator
both tend to $1$ and the denominator is never zero. In particular,
it exceeds $1/2$ eventually. Hence the displayed ratio is eventually
at least $(n+1)/2$, and therefore tends to $+\infty$.
The divergent part of the Ratio Test applies, since the ratio
limit inferior is greater than $1$.

We can also see why the term condition fails. There is $N$ after
which every ratio is at least $2$. Induction then gives
$a_{N+k}\ge2^k a_N$, with $a_N>0$, so the terms tend to
$+\infty$ rather than to zero. The positive partial sums are
unbounded, and the series diverges to $+\infty$.

\step{(e) The series with term $\cos^2 n/n^2$.}
The trigonometric bound $-1\le\cos n\le1$ gives
$0\le\cos^2 n\le1$. Since $n^2>0$, division preserves order:
\[
 0\le\frac{\cos^2 n}{n^2}\le\frac1{n^2}.
\]
The comparison series $\sum 1/n^2$ converges by the $p$-series
criterion with $p=2>1$. The Comparison Test therefore proves
convergence of the given nonnegative series.

\step{(f) The series with term $1/\log n$, beginning at $n=2$.}
For $n\ge2$, the logarithm is positive. The elementary logarithm
bound is
\[
 0<\log n=\int_1^n\frac{dt}{t}\le\int_1^n1\,dt=n-1<n.
\]
The integral comparison here only justifies the inequality for one
logarithm; no integral test for a series is being invoked.
Taking reciprocals of positive quantities reverses the inequality,
so $1/\log n\ge1/n$. The harmonic series diverges to $+\infty$,
and deleting its first term leaves divergent partial sums.
Thus the Comparison Test proves divergence to $+\infty$.

\answer{Parts (a), (b), (c), and (e) converge absolutely. Parts (d)
and (f) diverge to $+\infty$.}

\exercise{14.2}{Seven further convergence tests}
\goal Classify the seven series, explaining the actual numerator
and denominator in each term before applying a test.

\step{(a) The series with term $(n-1)/n^2$.}
All terms are nonnegative. For $n\ge2$, the inequality
$n-1\ge n/2$ follows by subtracting $n/2$ from both sides:
$n/2-1\ge0$. Therefore
\[
 \frac{n-1}{n^2}\ge\frac{n/2}{n^2}=\frac1{2n}\qquad(n\ge2).
\]
The series $\sum_{n=2}^{\infty}1/(2n)$ is one half of a tail of
the harmonic series, so it diverges to $+\infty$. The Comparison
Test proves divergence to $+\infty$ of the requested series.
The first term is zero and has no effect on this lower bound for
the later partial sums.

\step{(b) The series with term $(-1)^n$.}
Every term has absolute value $1$, so the terms do not tend to zero.
Corollary~14.5 therefore rules out convergence. Directly, the partial
sums are $-1,0,-1,0,\ldots$: an odd number of terms gives $-1$,
and each pair $(-1)+1$ sums to zero. Thus the partial sums have
two different subsequential limits and do not tend to either infinity.

\step{(c) The series with term $3n/n^3$.}
Since $n\ne0$, cancellation gives
$3n/n^3=3/n^2$. The $p$-series $\sum1/n^2$ converges, and
multiplication of all partial sums by the constant $3$ preserves
convergence by the constant multiple limit law. Hence this positive
series converges absolutely.

\step{(d) The series with term $n^3/3^n$.}
For $a_n=n^3/3^n>0$, compute
\[
 \frac{a_{n+1}}{a_n}
 =\frac13\left(1+\frac1n\right)^3
 \longrightarrow\frac13<1.
\]
All ratios are defined; the Ratio Test gives absolute convergence.

\step{(e) The series with term $n^2/n!$.}
Put $a_n=n^2/n!>0$. Cancel $n!$ and one factor of $n+1$:
\[
 \frac{a_{n+1}}{a_n}
 =\frac{(n+1)^2}{(n+1)!}\frac{n!}{n^2}
 =\frac{n+1}{n^2}
 =\frac1n+\frac1{n^2}\longrightarrow0<1.
\]
The Ratio Test proves absolute convergence.

\step{(f) The series with term $1/n^n$.}
The terms are positive and their positive $n$th roots satisfy
\[
 \left(\frac1{n^n}\right)^{1/n}=\frac1n\longrightarrow0.
\]
Thus the limit superior of the roots is $0<1$. The convergent
part of the Root Test proves absolute convergence.

\step{(g) The series with term $n^n/2^n$.}
Here the positive $n$th root is
\[
 \left(\frac{n^n}{2^n}\right)^{1/n}=\frac n2\longrightarrow+\infty.
\]
The root limit superior is greater than $1$, so the Root Test
proves divergence. More explicitly, for $n\ge4$ one has
$n/2\ge2$, so $(n/2)^n\ge2^n$. These positive terms tend to
$+\infty$; hence the partial sums do too.

\answer{Parts (c), (d), (e), and (f) converge absolutely.
Parts (a) and (g) diverge to $+\infty$. Part (b) diverges with
bounded, oscillating partial sums.}

\exercise{14.3}{Factorials, geometric comparisons, and a periodic term}
\goal Classify each of the six series and justify convergence or
failure of the necessary term condition.

\step{(a) The series with term $1/\sqrt{n!}$.}
Set $a_n=1/\sqrt{n!}>0$. The factorial is inside the square root.
Using $(n+1)!=(n+1)n!$ and positivity,
\[
 \frac{a_{n+1}}{a_n}
 =\frac{\sqrt{n!}}{\sqrt{(n+1)!}}
 =\frac1{\sqrt{n+1}}\longrightarrow0.
\]
The Ratio Test applies and gives absolute convergence.

\step{(b) The series with term $(2+\cos n)/3^n$.}
Since $-1\le\cos n\le1$, adding $2$ gives
$1\le2+\cos n\le3$. Consequently
\[
 0\le\frac{2+\cos n}{3^n}\le\frac3{3^n}.
\]
The comparison series is geometric, with positive constant factor
$3$ and ratio $1/3<1$. It converges, so the Comparison Test
proves absolute convergence of the given nonnegative series.

\step{(c) The series with term $1/(2^n+n)$.}
We have $2^n+n\ge2^n>0$. Reciprocation yields
\[
 0<\frac1{2^n+n}\le\frac1{2^n}.
\]
The upper comparison is a geometric series of ratio $1/2$,
so the requested series converges absolutely.

\step{(d) The series with term $(1/2)^n(50+2/n)$.}
For $n\ge1$, the inequality $0<2/n\le2$ implies
$0<50+2/n\le52$. Multiply by the positive number $2^{-n}$:
\[
 0<\left(\frac12\right)^n\left(50+\frac2n\right)
 \le\frac{52}{2^n}.
\]
The series on the right is a constant multiple of a convergent
geometric series. The Comparison Test therefore gives absolute
convergence. In particular, the expression in parentheses is
$50+2/n$, with $n$ in the denominator of its second term.

\step{(e) The series with term $\sin(n\pi/9)$.}
Consider the strictly increasing indices $n_k=18k+3$,
$k\in\N$. The corresponding terms are
\[
 \sin\left(\frac{(18k+3)\pi}{9}\right)
 =\sin(2k\pi+\pi/3)=\sin(\pi/3)=\frac{\sqrt3}{2}.
\]
This nonzero constant subsequence prevents the whole term sequence
from tending to zero: every subsequence of a sequence tending to
zero would also tend to zero. Thus Corollary~14.5 proves that the
series diverges. A term test does not require the terms to be
nonnegative or the partial sums to tend to $+\infty$.

\step{(f) The series with term $100^n/n!$.}
For $a_n=100^n/n!>0$, factorial cancellation gives
\[
 \frac{a_{n+1}}{a_n}=\frac{100}{n+1}\longrightarrow0<1.
\]
The Ratio Test gives absolute convergence, despite the relatively
large fixed base $100$ in the numerator.

\answer{Every series in this exercise except (e) converges
absolutely. Part (e) diverges because its terms do not tend to zero.}

\exercise{14.4}{Oscillating denominators, telescoping roots, and factorials}
\goal Classify the three series, keeping track of the starting
index and the placement of the powers.

\step{(a) The series with term $1/[n+(-1)^n]^2$, for $n\ge2$.}
For these indices,
\[
 n+(-1)^n\ge n-1\ge\frac n2>0.
\]
The middle inequality is equivalent to $n\ge2$. Because both
sides are positive, squaring preserves order and taking reciprocals
reverses it. We obtain
\[
 0<\frac1{[n+(-1)^n]^2}\le\frac4{n^2}\qquad(n\ge2).
\]
The comparison series converges by the $p$-series criterion with
$p=2$ and the constant multiple limit law. Hence the given series
converges absolutely. Beginning at $n=2$ also avoids the zero
denominator that would occur at $n=1$.

\step{(b) The series with term $\sqrt{n+1}-\sqrt n$.}
Each term is positive. Its $N$th partial sum telescopes:
\[
 \begin{aligned}
 \sum_{n=1}^{N}(\sqrt{n+1}-\sqrt n)
 &=(\sqrt2-1)+(\sqrt3-\sqrt2)+\cdots
     +(\sqrt{N+1}-\sqrt N)\\
 &=\sqrt{N+1}-1.
 \end{aligned}
\]
All intermediate square roots cancel exactly in this finite sum.
For a given $M>0$, any $N>(M+1)^2$ makes the last expression
greater than $M$. Thus the partial sums tend to $+\infty$ and
the series diverges to $+\infty$. Individual terms tending to
zero would not change this conclusion about the partial sums.

\step{(c) The series with term $n!/n^n$.}
Let $a_n=n!/n^n>0$. The ratio simplifies as follows:
\[
 \begin{aligned}
 \frac{a_{n+1}}{a_n}
 &=\frac{(n+1)!}{(n+1)^{n+1}}\frac{n^n}{n!}
 =\frac{(n+1)n^n}{(n+1)^{n+1}}\\
 &=\left(\frac{n}{n+1}\right)^n
 =\frac1{(1+1/n)^n}.
 \end{aligned}
\]
Exercise~12.14 justified $(1+1/n)^n\to e$ with finite $e\ge2$.
The reciprocal limit law therefore gives ratio limit $1/e$,
which lies in $(0,1)$. The Ratio Test proves absolute convergence.
\answer{Parts (a) and (c) converge absolutely. Part (b) diverges
to $+\infty$.}
